% New English translation from checked cote 25, site 103 (PDF page 104).
% AI-translated by GPT-6 Astra (OpenAI, ChatGPT Pro mode).
\scanpage{25/103}
\mspage{\hand{6}}
% NOTE: the reviewed source records a long diagonal pen line; its function does not establish cancellation of the entire argument.

and a homomorphism $f^{\hand{*}}_{\hand{*}}(\underline{L}) \hand{\longleftarrow} \underline{E}'$ (such that the \unsure{associated|associated} homomorphism
$f^{*}(\underline{E}') \hand{\to} \underline{L}$ is surjective), modulo isomorphism in $\underline{L}$. We will say that
$g\colon P \hand{\to} P'$ is \uline{of degree} n if $\underline{L} = g^{*}(\underline{O}_{P\unsure{(|'}}(1))$ is of degree n. It evidently
suffices to know how to determine the homomorphisms of degree n for every given
n. Note that if g is of degree n,\var{\inserted[typed]{and P has fibers of dimension $\hand{\geq} 1$}}{ and P has fibers of dimension $\hand{\geq} 1$}{insertion}{ega25_25_103-v0001} we will necessarily have $n \hand{\geq} 0$
(since over a field k, \var{\deleted{$\underline{O}_P(n)$}}{}{deletion}{ega25_25_103-v0002} if dim P $\hand{\geq} 1$, $\underline{O}_P(n)$ is generated by its
sections only if $n \hand{\geq} 0$); of course, we can \var{\deleted{also}}{}{deletion}{ega25_25_103-v0003} restrict to the case
where P has its fibers of dim $\hand{\geq} 1$ (proceeding as in 1.2.). Since we have
\[ f_{*}(f^{*}(\underline{M})(n)) = \underline{M} \otimes f_{*}(\underline{O}_P(n)) = \underline{M} \otimes \underline{\mathrm{Sym}}^{n}(\underline{E}) \hand{,} \]
we see that the determination of \unsure{S|g}-morphisms $g\colon P \hand{\to} P'$ reduces
to the determination of pairs $(\underline{M},u)$ up to isom, where $\underline{M}$ is an invertible
Module on S and $u\colon \underline{E}' \hand{\to} \underline{M}\otimes\underline{\mathrm{Sym}}^{n}(\underline{E})$ a homomorphism such that
the corresponding homomorphism $f^{*}(\underline{E}') \hand{\to} f^{*}(\underline{M})(n)$ is an epimorphism.
Thus g determines a first invariant of a global nature on S, namely
$\mathrm{cl}(\underline{M}) \hand{\in} \mathrm{Pic}(S)$, and this invariant being fixed by a chosen $\underline{M}$, the corresponding g's
correspond to a certain subset of the quotient set
$\mathrm{Hom}(\underline{E}',\underline{M}\otimes\underline{\mathrm{Sym}}^{n}(\underline{E}))/\mathrm{I}\,(S,\underline{O}_S)^{\hand{*}}$, passage to the quotient by the group
$\mathrm{I}\,(S,\underline{O}_S)^{\hand{*}}$ corresponding to the ``modulo isomorphism'' in the description
of the $g\colon P \hand{\to} P'$ via invertible Modules (NB the endomorphisms resp. automorphisms
of an invertible $\underline{L}$ on a projective bundle $\underline{P}$, corresponding to
sections resp. invertible\hand{/}sections of $\underline{O}_P$ on P, or again sections
resp invertible sections of $\underline{O}_S \hand{\xrightarrow{\sim}} f_{\hand{*}}(\underline{O}_P)$ on S).

Special cases: \hand{1)} n=0, then we must take the homomorphisms
$\underline{E}' \hand{\to} \underline{M}$ that are \var{\deleted{everywhere nonzero}}{}{deletion}{ega25_25_103-v0005}\var{\inserted{\uline{surjective} i.e. everywhere nonzero,}}{\uline{surjective} i.e. everywhere nonzero,}{insertion}{ega25_25_103-v0006}, modulo isomorphism in $\underline{M}$: we find
exactly the morphisms $g\colon P \hand{\to} P'$ of the form hf, where h is a
section of P' over S (namely the one determined by the invertible quotient
$\underline{M}$ of $\underline{E}'$). Thus the \var{\inserted[typed]{S-}}{S-}{insertion}{ega25_25_103-v0007}\uline{morphisms of degree 0} \uline{from} P \uline{to} P' \uline{are the morphis-}\var{\margin{\ensuremath{\backslash}}}{\placedmargin{\ensuremath{\backslash}}}{marginal note placed}{ega25_25_103-v0008}%
\tnote{25-103-1}{The uncertain spellings associé or assooié give the same English word. The canceled French t or x after au has no English grammatical effect; the source doubt is retained here as \unsure{t|x}. The two star positions in the opening expression, the parenthesis or prime-like mark by P, and the roman I in the quotient notation follow the source. The doubled comma after the inserted everywhere nonzero phrase remains.}

\var{\margin{103}}{\placedmargin{103}}{marginal note placed}{ega25_25_103-v0009}
